% ORDER3EXACT1
\documentclass[11pt]{article}
\usepackage[margin=.88in]{geometry}
\usepackage{amsmath,amssymb,amsthm}
\newtheorem{proposition}{Proposition}
\begin{document}
\begin{center}
{\Large ORDER3EXACT1: Exact third-order census after pair matching}\\[4pt]
\textbf{Date: 2026-10-04}\\
\textbf{Status: CANDIDATE lemmas + COMPUTATIONAL pilot}\\
\textbf{P versus NP status: OPEN}
\end{center}

\section{Motivation}
PAIR-CAMOUFLAGE2 and PAIRMATCH3 found a stable finite second-order signal in rare exact-layer UNSAT formulas at density
\[
t=3.5q,
\]
then substantially neutralized it by pair matching.  This run asks whether a complete, rather than sampled, three-clause orbit census reveals a reproducible third-order signal.

\section{Low-degree interpretation on the slice}
The exact Hamming layer is the Johnson slice. Filmus and collaborators develop harmonic analysis on this domain, including an orthogonal basis and invariance principles for low-degree functions. Filmus--Ihringer and Filmus further show that Boolean constant-degree functions on sufficiently large slices are juntas.

For two probability distributions \(\mu_+,\mu_-\) on the same exact slice, suppose
\[
\mathbb E_{\mu_+}\prod_{i\in S}z_i
=
\mathbb E_{\mu_-}\prod_{i\in S}z_i
\qquad
\forall |S|\le d.
\]

\begin{proposition}[MOMENT-CAMOUFLAGE]
Under the preceding hypothesis, every real multilinear polynomial of degree at most \(d\) has identical expectation under \(\mu_+\) and \(\mu_-\).
\end{proposition}
\begin{proof}
Expand the polynomial in the squarefree monomial basis. On \(\{0,1\}^M\), every monomial reduces to \(\prod_{i\in S}z_i\). Equality of all corresponding moments through degree \(d\) gives equality by linearity.
\end{proof}

Thus the empirical hierarchy
\[
\text{order 1}\to\text{order 2}\to\text{order 3}\to\cdots
\]
has an exact algebraic meaning: matching all moments through order \(d\) exactly fools every degree-\(d\) polynomial test. This does not imply that general polynomial-size monotone circuits have bounded harmonic degree.

\section{Exact three-clause census}
The complete three-clause orbit canonicalization from PAIRMATCH3 was used. For every retained formula, all
\[
\binom t3
\]
clause triples were enumerated.

The intended larger \(q=20,24\) run exceeded the available execution window when combined with rare-event sampling. The exact census was therefore completed as a pilot at \(q=16,18\), with tiny hostile checks at \(q=20,24\). No asymptotic extrapolation is made.

Pair matching used the dominant six standardized pair coordinates with caliper
\[
d_{\rm pair}\le1.0.
\]

\[
\begin{array}{c|c|c|c|c|c|c}
q&N&U_{hold}&U_{keep}&\#\text{orbits seen}&
\text{score BA}&\text{logistic BA}\\
\hline
16&1000&28&14&82&0.643&0.571\\
18&1200&15&12&80&0.417&0.500\\
20&700&4&3&62&0.500&0.500\\
24&900&4&4&69&0.500&0.500
\end{array}
\]

The score classifier selected the twelve strongest exact triple-orbit deviations on the discovery half and evaluated them on disjoint matched pairs. The logistic classifier used those same selected exact orbit counts.

The isolated \(q=16\) score value \(0.643\) does not replicate. The other tested values are at or below chance, although \(q=20,24\) are severely underpowered.

\section{Aggressive matched-core search}
A separate pair-matched run performed many randomized clause-deletion orders and retained the smallest deletion-minimal core found.

At \(q=16\), 100 deletion orders per retained UNSAT formula gave
\[
\overline n_{core}=13.67,
\qquad
n_{min}=11,
\]
\[
\overline m_{core}=24.27,
\qquad
m_{min}=18.
\]

At \(q=18\), 50 deletion orders per retained UNSAT formula gave
\[
\overline n_{core}=15.63,
\qquad
n_{min}=14,
\]
\[
\overline m_{core}=27.25,
\qquad
m_{min}=24.
\]

These are the smallest cores found by heuristic multi-start deletion, not proofs of minimum MUS size.

\section{Interpretation}
The finite hierarchy is now
\[
\boxed{\text{order 1: exactly matched by symmetry}},
\]
\[
\boxed{\text{order 2: strong raw signal, largely removed by matching}},
\]
\[
\boxed{\text{order 3: no replicated exact-census signal found yet}}.
\]

This is the first stage where a complete fixed-order orbit census, rather than a sampled feature library, fails to show a stable signal in the adequately populated pilot sizes. Positive counts remain small, so this is not a proof of asymptotic degree-three camouflage.

\section{Literature connection}
Filmus--Mossel prove an invariance principle for low-degree functions on the slice, and Filmus--Ihringer / Filmus prove junta theorems for Boolean constant-degree slice functions. These results make low-degree harmonic structure a natural hierarchy on our exact-weight domain.

They do not provide the missing implication
\[
\text{small general monotone circuit}
\Longrightarrow
\text{low harmonic degree}.
\]
An approximation theorem of that strength would itself be a major lower-bound result. The robust-sunflower / approximation route remains the more plausible source of a circuit-to-local-structure bridge.

\section{Next targets}
\begin{enumerate}
\item \textbf{ORDER3EXACT2:} increase the rare-UNSAT sample at \(q=20,24\) until the exact triple census has several dozen matched positive evaluation formulas.
\item \textbf{MOMENT-BALANCE3:} replace nearest-neighbor matching by explicit weighting/rejection so all degree-one and degree-two orbit moments match on training data, then test the complete degree-three vector on held-out formulas.
\item \textbf{MATCHED-MUF2:} use exact or branch-and-bound MUS search at smaller \(q\) to distinguish failure to find a small core from proof that none exists.
\item \textbf{SAT-SUNFLOWER2:} only after exact third-order tests remain quiet at larger positive sample sizes.
\end{enumerate}

\section{Status}
MOMENT-CAMOUFLAGE and the exact triple-orbit representation are \textbf{CANDIDATE with complete elementary proofs}. The census and core results are \textbf{COMPUTATIONAL}. Asymptotic degree-three camouflage, minimum-core lower bounds, and the general monotone circuit lower bound remain \textbf{OPEN}. Therefore \(P\ne NP\) remains OPEN.

\begin{thebibliography}{4}
\bibitem{FM2016} Y. Filmus and E. Mossel, Harmonicity and Invariance on Slices of the Boolean Cube, CCC 2016.
\bibitem{F2016} Y. Filmus, An Orthogonal Basis for Functions over a Slice of the Boolean Hypercube, Electronic Journal of Combinatorics 23(1), 2016.
\bibitem{FI2019} Y. Filmus and F. Ihringer, Boolean Constant Degree Functions on the Slice Are Juntas, Discrete Mathematics 342(12), 2019.
\bibitem{F2023} Y. Filmus, Junta Threshold for Low Degree Boolean Functions on the Slice, Electronic Journal of Combinatorics 30(1), 2023.
\end{thebibliography}
\end{document}
